% ch05_u01 —— 原书 p.192–p.200（第 5 章开头：5.1 全部；5.2 全部；5.3 自 5.3.1 起；p.200 末句中断于「So by」）
\bdchap{ch:hecke}{Hecke 特征 D-模}
\SourcePage{192}

\bdsec{5.1. D-模的构造}

\bdp{5.1.1} 本小节在 $\operatorname{Bun}_G$ 上构造一族 $\mathcal{D}$-模，以 $\mathcal{O}\mathfrak{p}_{{}^{L}G}(X)$ 为参数；后者即 $X$ 上 ${}^{L}G$-oper 的栈。

记 $Z$ 为 ${}^{L}G$ 的中心。按 3.4.3 的公式 (57)，须对 $\mathcal{L}\in Z\operatorname{tors}_\theta(X)$ 配置 $\operatorname{Bun}_G$ 上一族 $\mathcal{D}$-模，以 $\mathcal{O}\mathfrak{p}_{{}^{L}\mathfrak{g}}(X)$ 为参数。在 4.4.3 中我们定义了 $\lambda_{\mathcal{L}}\in\mu_\infty\operatorname{tors}_\theta(\operatorname{Bun}_G)$。$\lambda_{\mathcal{L}}$ 是 $\operatorname{Bun}_G$ 上的线丛，带有同构 $\lambda_{\mathcal{L}}^{\otimes 2n}\xrightarrow{\sim}(\omega^{\sharp}_{\operatorname{Bun}_G})^{\otimes n}$，其中 $n\ne 0$（见 4.0.1）。故 $\lambda_{\mathcal{L}}$ 是 $\mathcal{D}'$-模。因此 $M_{\mathcal{L}}:=\lambda_{\mathcal{L}}^{-1}\otimes_{\mathcal{O}_{\operatorname{Bun}_G}}\mathcal{D}'$ 是 $\operatorname{Bun}_G$ 上的左 $\mathcal{D}$-模。按 3.3.2 与 2.7.4，有典范的代数态射 $h_X\varphi_X:A_{{}^{L}\mathfrak{g}}(X)\to\Gamma(\operatorname{Bun}_G,\mathcal{D}')$。于是 $\Gamma(\operatorname{Bun}_G,\mathcal{D}')$ 在 $\mathcal{D}'$ 上的右作用，给出 $M_{\mathcal{L}}$ 上的 $A_{{}^{L}\mathfrak{g}}(X)$-模结构。因此可以把 $M_{\mathcal{L}}$ 看作 $\operatorname{Bun}_G$ 上一族左 $\mathcal{D}$-模，以 $\operatorname{Spec}A_{{}^{L}\mathfrak{g}}(X)=\mathcal{O}\mathfrak{p}_{{}^{L}\mathfrak{g}}(X)$ 为参数。
% [Tutor Note: 原书 p.192 把复合态射印作 $h_X\varphi_X$（$h_X$ 与 $\varphi_X$ 之间无复合记号），底本如此，保留原样。]

这样，我们在 $\operatorname{Bun}_G$ 上构造了一族左 $\mathcal{D}$-模，以 $\mathcal{O}\mathfrak{p}_{{}^{L}G}(X)$ 为参数。对 ${}^{L}G$-oper $\mathfrak{F}$，相应的 $\mathcal{D}$-模 $M_{\mathfrak{F}}$ 是 $M_{\mathcal{L}}/\mathfrak{m}_{\mathfrak{F}}M_{\mathcal{L}}=\lambda_{\mathcal{L}}^{-1}\otimes\mathcal{D}'/\mathcal{D}'\mathfrak{m}_{\mathfrak{F}}$，其中 $\mathcal{L}$ 是 $\mathfrak{F}$ 在 $Z\operatorname{tors}_\theta(X)$ 中的像，而 $\mathfrak{m}_{\mathfrak{F}}\subset A_{{}^{L}\mathfrak{g}}(X)$ 是对应于 $\mathfrak{F}$ 的 ${}^{L}\mathfrak{g}$-oper 的极大理想。

\bdp{5.1.2} \emph{命题.}
\begin{enumerate}[label=(\roman*),leftmargin=2.2em]
\item 对每个 $\mathcal{L}\in Z\operatorname{tors}_\theta(X)$，$M_{\mathcal{L}}$ 在 $A_{{}^{L}\mathfrak{g}}(X)$ 上平坦。
\item 对每个 ${}^{L}G$-oper $\mathfrak{F}$，$\mathcal{D}$-模 $M_{\mathfrak{F}}$ 完整；其奇异支集作为闭链，与 Hitchin 纤维化的零纤维重合。
\end{enumerate}

\emph{证明.} 按 2.2.4(iii)，$\operatorname{gr}\mathcal{D}'$ 在 $\operatorname{gr}A_{{}^{L}\mathfrak{g}}(X)$ 上平坦\footnote{这意指：若 $f:S\to\operatorname{Bun}_G$ 光滑且 $S$ 仿射，则 $\Gamma(S,f^*\operatorname{gr}\mathcal{D}')$ 是自由 $\operatorname{gr}A_{{}^{L}\mathfrak{g}}(X)$-模（在 $A_0=\mathbb{C}$ 的 $\mathbb{Z}_+$-分次环 $A$ 上，平坦的 $\mathbb{Z}_+$-分次模是自由的）。}。故 $\mathcal{D}'$ 在 $A_{{}^{L}\mathfrak{g}}(X)$ 上平坦。这蕴含 i)，以及对任意理想 $I\subset A_{{}^{L}\mathfrak{g}}(X)$ 的等式 $\operatorname{gr}(\mathcal{D}'/\mathcal{D}'I)=\operatorname{gr}\mathcal{D}'/(\operatorname{gr}\mathcal{D}'\cdot\operatorname{gr}I)$。若 $I$ 极大，即得 ii)。\hfill$\Box$

\SourcePage{193}

\bdsec{5.2. 主要定理 I：导引}

\bdp{5.2.1} 整体主定理 5.2.6 断言：$\mathcal{D}$-模 $M_{\mathfrak{F}}$ 是\emph{Hecke 函子}的特征模。为定义 Hecke 函子，我们引入\emph{大 Hecke 栈} $\mathcal{H}ecke$。$S$-点构成的群胚 $\mathcal{H}ecke(S)$ 由四元组 $(\mathcal{F}_1,\mathcal{F}_2,x,\alpha)$ 组成，其中 $\mathcal{F}_1,\mathcal{F}_2$ 是 $X\times S$ 上的 $G$-torsor，$x\in X(S)$，而 $\alpha:\mathcal{F}_1|_U\xrightarrow{\sim}\mathcal{F}_2|_U$ 是 $x$ 的图像之补集 $U$ 上的同构。有显然的投影 $p_{1,2,X}=(p_1,p_2,p_X):\mathcal{H}ecke\to\operatorname{Bun}_G\times\operatorname{Bun}_G\times X$。

栈 $\mathcal{H}ecke$ 是归纳代数的，投影 $p_i$、$p_{i,X}$ 是归纳固有的。确切地说，存在闭代数子栈的递增族 $\mathcal{H}ecke_1\subset\mathcal{H}ecke_2\subset\cdots\subset\mathcal{H}ecke$，使得 $\mathcal{H}ecke=\bigcup\mathcal{H}ecke_a$，且 $p_i:\mathcal{H}ecke_a\to\operatorname{Bun}_G$ 与 $p_{i,X}:\mathcal{H}ecke_a\to\operatorname{Bun}_G\times X$ 都是固有态射。

\bdp{5.2.2} \emph{注.} (i) 诸 $\alpha$ 的复合使 $\mathcal{H}ecke$ 成为 $\operatorname{Bun}_G$ 上的 $X$-群胚族。

(ii) $\mathcal{H}ecke$ 是 $\operatorname{Bun}_G\times X$ 上扭仿射 Grassmann 流形的族。确切地说，对 $(\mathcal{F}_2,x)\in\operatorname{Bun}_G\times X$，纤维 $\mathcal{H}ecke_{(\mathcal{F}_2,x)}:=p_{2,X}^{-1}(\mathcal{F}_2,x)$ 典范同构于仿射 Grassmann 流形 $\mathcal{GR}_x:=G(K_x)/G(O_x)$ 经 $G(O_x)$-torsor $\mathcal{F}_2(O_x)$ 所作的扭（就左 $G(O_x)$-作用而言）。当 $\mathcal{F}_2$ 是平凡丛时，我们在 4.5.2 中描述了这个同构。一般情形的构造类似：固定 $\gamma_2\in\mathcal{F}_2(O_x)$，把 $(\mathcal{F}_1,\mathcal{F}_2,x,\alpha)$ 对应到 $\gamma_2/\alpha(\gamma_1)$ 在 $G(K_x)/G(O_x)$ 中的像；这里 $\gamma_1$ 是 $\mathcal{F}_1(O_x)$ 的任一元素，$\gamma_2/\alpha(\gamma_1)$ 表示满足 $g\alpha(\gamma_1)=\gamma_2$ 的元素 $g\in G(K_x)$。按 2.3.4，态射 $\mathcal{H}ecke_{(\mathcal{F}_2,x)}\to G(K_x)/G(O_x)$ 是同构。

\bdp{5.2.3} 态射 $\nu:\mathbb{G}_m\to G$ 的共轭类集，可典范地等同于 ${}^{L}G$ 的支配权集 $P_+({}^{L}G)$。回忆 $\mathcal{GR}_x=G(K_x)/G(O_x)$ 中的 $G(O_x)$-轨道由 $\chi\in P_+({}^{L}G)$ 标号。按定义，$\operatorname{Orb}_\chi$ 是 $\nu(t_x)\in G(K_x)$ 在 $\mathcal{GR}_x$ 中的像的 $G(O_x)$-轨道，其中 $\nu:\mathbb{G}_m\to G$ 的类为 $\chi$，$t_x\in O_x$ 是素元。

\SourcePage{194}

按 5.2.2(ii)，$\mathcal{GR}_x$ 按 $\operatorname{Orb}_\chi$ 所作的分层，给出栈 $\mathcal{H}ecke$ 按子栈 $\mathcal{H}ecke_\chi$ 的分层，$\chi\in P_+({}^{L}G)$。$\mathcal{H}ecke_\chi$ 的 $\mathbb{C}$-点是四元组 $(\mathcal{F}_1,\mathcal{F}_2,x,\alpha)$，使得对某组 $\gamma_i\in\mathcal{F}_i(O_x)$ 与 $x$ 处某个形式参数 $t_x$，有 $\gamma_2=\nu(t_x)\alpha(\gamma_1)$，其中 $\nu:\mathbb{G}_m\to G$ 的类为 $\chi$。对合 $(\mathcal{F}_1,\mathcal{F}_2,x,\alpha)\mapsto(\mathcal{F}_2,\mathcal{F}_1,x,\alpha^{-1})$ 把 $\mathcal{H}ecke_\chi$ 与 $\mathcal{H}ecke_{\chi^\circ}$ 等同，其中 $\chi^\circ$ 是对偶权。于是 $p_{2,X}:\mathcal{H}ecke_\chi\to\operatorname{Bun}_G\times X$ 的纤维是 $\operatorname{Orb}_\chi$ 的扭形式，而 $p_{1,X}:\mathcal{H}ecke_\chi\to\operatorname{Bun}_G\times X$ 的纤维是 $\operatorname{Orb}_{\chi^\circ}$ 的扭形式。

对每个 $\chi$，栈 $\mathcal{H}ecke_\chi$ 在 $\operatorname{Bun}_G\times X$ 上光滑。通常其闭包 $\overline{\mathcal{H}ecke_\chi}$ 不光滑。
% [Tutor Note: 原书 p.194 此处既写出「closure」一词，又用上划线记号 $\overline{\mathcal{H}ecke_\chi}$，语义重复；底本如此，照录。]

\emph{注.} (i) 按 4.5.12，$\overline{\mathcal{H}ecke_\chi}$ 是诸分层 $\mathcal{H}ecke_{\chi'}$（$\chi'\le\chi$）之并。

(ii) 若 $G=GL(n)$，则我们的分层标号与「自然」标号一致。确切地说，设 $V_1,V_2$ 是对应于 $\mathfrak{F}_1,\mathfrak{F}_2$ 的向量丛。则 $\mathcal{H}ecke_\chi$ 由一切这样的组 $(V_1,V_2,x,\alpha)$ 组成：在 $x$ 的形式邻域上适当选取 $V_i$ 的基之后，$\alpha$ 的矩阵等于 $t_x^\chi$。

\bdp{5.2.4} 下面定义 Hecke 函子 $T_\chi^i:\mathcal{M}(\operatorname{Bun}_G)\to\mathcal{M}(\operatorname{Bun}_G\times X)$，其中 $\mathcal{M}$ 表示 $\mathcal{D}$-模范畴，$\chi\in P_+({}^{L}G)$，$i\in\mathbb{Z}$。

对 $\chi\in P_+({}^{L}G)$、$M\in\mathcal{M}(\operatorname{Bun}_G)$，记 $p_{1\chi}^\star M$ 为下述对象：$M$ 经光滑投影 $p_{1\chi}:\mathcal{H}ecke_\chi\to\operatorname{Bun}_G$ 的拉回，在 $\overline{\mathcal{H}ecke_\chi}$ 上的极小（$=$ Goresky--MacPherson）扩张；这里 $p_{1\chi}:=p_1|_{\mathcal{H}ecke_\chi}$。注意纤维化 $p_{1X}:\overline{\mathcal{H}ecke_\chi}\to\operatorname{Bun}_G\times X$ 局部平凡（见 5.2.2(ii)、5.2.3）。因此选定局部平凡化后，$p_{1\chi}^\star M$（局部地）等同于 $M$ 与仿射 Grassmann 流形上相应 $G(O)$-轨道\footnote{该轨道是 $\operatorname{Orb}_{\chi^\circ}$，其中 $\chi^\circ$ 是对偶权，见 5.2.3。}之闭包上的「相交上同调」$\mathcal{D}$-模的外张量积。

Hecke 函子 $T_\chi^i:\mathcal{M}(\operatorname{Bun}_G)\to\mathcal{M}(\operatorname{Bun}_G\times X)$ 定义为
\begin{equation}\tag{254}
T_\chi^i=H^i(p_{2,X})_*p_{1\chi}^\star
\end{equation}

\SourcePage{195}

其中 $H^i(p_{2,X})_*$ 是投影 $p_{2,X}:\overline{\mathcal{H}ecke_\chi}\to\operatorname{Bun}_G\times X$ 的上同调前推函子。

\emph{注.} 对局部有限型代数栈的可表拟紧态射 $f:\mathcal{X}\to\mathcal{Y}$，$H^if_*:\mathcal{M}(\mathcal{X})\to\mathcal{M}(\mathcal{Y})$ 的定义是清楚的。事实上，概形情形已有 $H^if_*$ 的定义，且已知 $H^if_*$ 与光滑基变换交换。

\bdp{5.2.5} 对 $\chi\in P_+({}^{L}G)$，记 $V^\chi$ 为最高权 $\chi$ 的不可约 ${}^{L}G$-模，并带一个标记的最高向量。若 $\mathfrak{F}$ 是 $X$ 上的 ${}^{L}G$-oper（更一般地，带联络的 ${}^{L}G$-丛），记 $V^\chi_{\mathfrak{F}}$ 为 $V^\chi$ 的 $\mathfrak{F}$-扭；它是 $X$ 上的光滑 $\mathcal{D}$-模。

\bdp{5.2.6} \emph{整体主定理.} 设 $\mathfrak{F}$ 是 $X$ 上的 ${}^{L}G$-oper，$M_{\mathfrak{F}}$ 是 5.1.1 中定义的 $\operatorname{Bun}_G$ 上的 $\mathcal{D}$-模。则当 $i\ne 0$ 时 $T_\chi^iM_{\mathfrak{F}}=0$，且有 $\operatorname{Bun}_G\times X$ 上 $\mathcal{D}$-模的典范同构
\begin{equation}\tag{255}
T_\chi^0M_{\mathfrak{F}}\xrightarrow{\sim}M_{\mathfrak{F}}\boxtimes V^\chi_{\mathfrak{F}}.
\end{equation}

同构 (255) 与 Hecke 对应的复合以及 $V^\chi$ 的张量积相容。确切陈述见 5.4.3。这一切意味着：$M_{\mathfrak{F}}$ 是以 $\mathfrak{F}$ 为特征值的\emph{Hecke 特征 $\mathcal{D}$-模}。

\bdp{5.2.7} Laumon 对 $X$ 上不可约局部系统 $E$ 定义了一个猜测中的「朗兰兹变换」$K_E$（见 [La87] 的 §5.3 与 4.3.3；当 $\operatorname{rank}E\le 2$ 时 $K_E$ 确实存在）。$K_E$ 是 $\operatorname{Bun}_{GL_n}$ 上的完整 $\mathcal{D}$-模，$n=\operatorname{rank}E$；至少当 $n=2$ 时，其奇异支集是 Hitchin 纤维化的零纤维（见 [La87] 的 §5.5）。此外 $K_E$ 有正则奇点，且它在 $\operatorname{Bun}_{GL_n}$ 每个连通分支上的限制不可约。若 $E$ 是 $SL_n$-局部系统，则 $K_E$ 位于 $\operatorname{Bun}_{PGL_n}$ 上。

考虑 5.1.2 与 5.2.6，自然可以猜测：当 $G=PGL_n$ 时，5.1.1 中的 $\mathcal{D}$-模 $M_{\mathfrak{F}}$ 等于 $K_{\mathfrak{F}}$（这方面的一些结果可见 [Fr]）。此外，还想弄清：\SourcePage{196}（对任意 $G$）$M_{\mathfrak{F}}$ 是否有正则奇点，它在 $\operatorname{Bun}_G$ 各连通分支上的限制是否不可约。

\bdp{5.2.8} 用 5.1.1 中的 $\mathcal{D}$-模 $M_{\mathcal{L}}$（$\mathcal{L}\in Z\operatorname{tors}_\theta(X)$）改写 5.2.6，既方便又重要。按 (57)，$\mathcal{L}\in Z\operatorname{tors}_\theta(X)$ 定义了 $X$ 上一族 ${}^{L}G$-oper $\mathfrak{F}_{\mathcal{L}}$，以 $\operatorname{Spec}A_{{}^{L}\mathfrak{g}}(X)$ 为参数。于是 $\mathfrak{F}_{\mathcal{L}}$ 是 $X\times\operatorname{Spec}A_{{}^{L}\mathfrak{g}}(X)$ 上的 ${}^{L}G$-torsor，带有沿 $X$ 的联络。对 $\chi\in P_+({}^{L}G)$，$V^\chi$ 的 $\mathfrak{F}_{\mathcal{L}}$-扭是 $X\times\operatorname{Spec}A_{{}^{L}\mathfrak{g}}(X)$ 上的向量丛，带有沿 $X$ 的联络。我们把它看作 $X$ 上的 $\mathcal{D}$-模 $V^\chi_{\mathcal{L}}$，并带有 $A_{{}^{L}\mathfrak{g}}(X)$ 的作用。

现在考虑 $\operatorname{Bun}_G$ 上的 $\mathcal{D}$-模 $M_{\mathcal{L}}$（见 5.1.1）；$A_{{}^{L}\mathfrak{g}}(X)$ 作用于它。容易看出（用 5.1.2(i)）：5.2.6 是下述定理的推论。

\bdp{5.2.9} \emph{定理.} 存在 $\operatorname{Bun}_G\times X$ 上 $\mathcal{D}$-模的典范同构
\begin{equation}\tag{256}
T_\chi^0M_{\mathcal{L}}\simeq M_{\mathcal{L}}\underset{A_{{}^{L}\mathfrak{g}}(X)}{\boxtimes}V^\chi_{\mathcal{L}},
\end{equation}
它与 $A_{{}^{L}\mathfrak{g}}(X)$ 的作用相容；并且当 $i\ne 0$ 时 $T_\chi^iM_{\mathcal{L}}=0$。

\bdp{5.2.10} 我们将从整体定理的局部版本推出它，现在就来说明这个局部版本。考虑仿射 Grassmann 流形 $\mathcal{GR}:=G(K)/G(O)$，其中 $O:=\mathbb{C}[[t]]$，$K=\mathbb{C}((t))$。这是归纳固有的归纳概形。于是有 $\mathcal{GR}$ 上 $\mathcal{D}$-模的「抽象」范畴 $\mathcal{M}(\mathcal{GR})$，它定义为 $\varinjlim\mathcal{M}(Y)$，其中 $Y$ 取遍 $\mathcal{GR}$ 的一切闭子概形。

我们无法把 $\mathcal{GR}$ 表示成光滑子概形递增序列之并。然而 $\mathcal{GR}$ 是\emph{形式光滑}的归纳概形。这允许把 $\mathcal{GR}$ 上的 $\mathcal{D}$-模当作「具体」对象来处理，就像 $\mathcal{GR}$ 是光滑有限维簇一样；也就是说，把它们等同于带某些附加结构的 $\mathcal{O}$-模层。确切地说，设 $\mathcal{GR}$ 上有 $\mathcal{O}$-模 $P$，使得 $P$ 的每个局部截面\SourcePage{197}都支在 $\mathcal{GR}$ 的某个子概形上。于是容易定义 $\operatorname{Der}\mathcal{O}_{\mathcal{GR}}$ 在 $P$ 上的连续\emph{右}作用。这样的 $P$ 带上这样的作用，就等同于 $\mathcal{GR}$ 上的 $\mathcal{D}$-模（我们还假定适当的拟凝聚性条件）。细节可见 ???。
% [Tutor Note: 原书 p.197「Details can be found in ???」中的交叉引用印作三个问号、无编号，底本如此，保留原样。]

\bdp{5.2.11} \emph{注.} 我们看到，在这个无穷维情形下，只有\emph{右} $\mathcal{D}$-模才真正作为层有意义。其原因相当有限维。事实上，设 $i:Y\hookrightarrow Z$ 是光滑流形的闭嵌入，$M$ 是 $Y$ 上的 $\mathcal{D}$-模；为把 $M$ 等同于 $i_*M$ 的子层，需要用到\emph{右} $\mathcal{D}$-模。

\bdp{5.2.12} 按 3.4.3，有群胚 $Z\operatorname{tors}_\theta(O)$，它是 $Z\operatorname{tors}_\theta(X)$ 的局部类比。选定 $\mathcal{L}\in Z\operatorname{tors}_\theta(O)$（本质上相当于选定 $\omega_O$ 的平方根）后，就定义了 $\mathcal{GR}$ 上「局部」Pfaffian 线丛 $\lambda_{\mathcal{L}}^{\operatorname{loc}}$（见 4.6）。$\mathfrak{g}\otimes K$ 经左无穷小平移在 $\mathcal{GR}$ 上的作用，提升为 2.5.1 的中心扩张 $\widetilde{\mathfrak{g}\otimes K}$ 在 $\lambda_{\mathcal{L}}^{\operatorname{loc}}$ 上的作用，使得 $\mathbf{1}\in\mathbb{C}\subset\widetilde{\mathfrak{g}\otimes K}$ 的作用为乘以 $-1$（见 4.6.5）。这给出反同态 $\bar U'\to\Gamma(\mathcal{GR},\mathcal{D}')$，其中 $\bar U'=\bar U'(\mathfrak{g}\otimes K)$ 是 2.9.4 中定义的完备化扭泛包络代数，而 $\Gamma(\mathcal{GR},\mathcal{D}')$ 是 $\mathcal{GR}$ 上 $\lambda_{\mathcal{L}}^{\operatorname{loc}}$-扭微分算子构成的环。因此，对 $\mathcal{GR}$ 上任一 $\mathcal{D}$-模 $M$，代数 $\bar U'$ 作用于 $M\lambda_{\mathcal{L}}^{-1}:=M\otimes_{\mathcal{O}_{\mathcal{GR}}}(\lambda_{\mathcal{L}}^{\operatorname{loc}})^{\otimes -1}$。故 $\Gamma(\mathcal{GR},M\lambda_{\mathcal{L}}^{-1})$ 是（左）$\bar U'$-模。

例如，考虑 $\mathcal{GR}$ 的基点处 $\delta$-函数构成的 $\mathcal{D}$-模 $I_1$。$\bar U'$-模 $\Gamma(\mathcal{GR},I_1\lambda_{\mathcal{L}}^{-1})$ 就是真空模 $Vac'$。

\bdp{5.2.13} 回忆（见 4.5.8）：$\mathcal{GR}$ 按 $G(O)$-轨道 $\operatorname{Orb}_\chi$ 分层，轨道由 $\chi\in P_+({}^{L}G)$ 标号。记 $I_\chi$ 为 $\mathcal{GR}$ 上对应于 $\overline{\operatorname{Orb}}_\chi$ 的不可约「相交上同调」$\mathcal{D}$-模。

以下是局部主定理的第一部分。

\bdp{5.2.14} \emph{定理.} $\bar U'$-模 $\Gamma(\mathcal{GR},I_\chi\lambda_{\mathcal{L}}^{-1})$ 同构于若干 $Vac'$ 的直和；并且当 $i>0$ 时 $H^i(\mathcal{GR},I_\chi\lambda_{\mathcal{L}}^{-1})=0$。

\SourcePage{198}

\emph{注.} 该定理意味着（见 5.4.8、5.4.10）：Harish-Chandra 模 $Vac'$ 是 7.8.2、7.14.1 中 Hecke 函子之 Harish-Chandra 版本的特征模。

\bdp{5.2.15} 群 $\operatorname{Aut}O$ 作用于 $\mathcal{GR}$，其李代数 $\operatorname{Der}O$ 的作用提升到 $\lambda_{\mathcal{L}}^{\operatorname{loc}}$ 上（见 4.6.7）。定理的第二部分描述 $\operatorname{Der}O$ 在 $\Gamma(\mathcal{GR},I_\chi\lambda_{\mathcal{L}}^{-1})$ 上的作用。

考虑 3.2.1 中局部 ${}^{L}\mathfrak{g}$-oper 的概形 $\mathcal{O}\mathfrak{p}_{{}^{L}\mathfrak{g}}(O)=\operatorname{Spec}A_{{}^{L}\mathfrak{g}}(O)$。把 $A_{{}^{L}\mathfrak{g}}(O)$ 简记为 $A$。与 5.2.8 一样，$\mathcal{L}$ 定义了 $\operatorname{Spec}O$ 上一族 ${}^{L}G$-oper，以 $\operatorname{Spec}A$ 为参数。这族 oper 定义了 $\operatorname{Spec}A$ 上的 ${}^{L}G$-torsor $\mathfrak{F}_A$，带有 $\operatorname{Der}O$ 的作用，且该作用与 $\operatorname{Der}O$ 在 $A$ 上的作用相容；见 3.5.4\footnote{在 3.5.4 中我们用的记号是 $\mathfrak{F}_G^0$ 而非 $\mathfrak{F}_A$，并考虑了 $\mathcal{L}$ 为 $\omega_O$ 的平方根这一「特殊」情形。}。${}^{L}G$-模 $V^\chi$ 的 $\mathfrak{F}_A$-扭是 $\operatorname{Spec}A$ 上的向量丛。记 $V^\chi_{\mathcal{L}A}$ 为其截面构成的 $A$-模；$\operatorname{Der}O$ 作用于它。

\bdp{5.2.16} \emph{定理.} 存在 $\bar U'$-模的典范同构
\begin{equation}\tag{257}
\Gamma(\mathcal{GR},I_\chi\lambda_{\mathcal{L}}^{-1})\xrightarrow{\sim}Vac'\otimes_AV^\chi_{\mathcal{L}A},
\end{equation}
它与 $\operatorname{Der}O$ 的作用相容。

这里我们用到 $Vac'$ 上的 $A$-模结构，它来自 Feigin--Frenkel 同构 (80)。

\bdp{5.2.17} 关于证明略说几句。整体定理由局部定理经一个容易的局部--整体论证得到，该论证与 2.8 中所用者类似。局部定理的证明基于下述两个关键结构之间的相互作用：
\begin{enumerate}[label=(\roman*),leftmargin=2.2em]
\item \emph{Satake 等价}（[Gi95]、[MV]）：${}^{L}G$ 的表示张量范畴，与由诸 $I_\chi$ 生成、带「卷积」张量结构的 $\mathcal{GR}$ 上 $\mathcal{D}$-模范畴之间的等价。
\item \emph{「重正化」包络代数} $U^\natural$。代数态射 $\bar U'\to\Gamma(\mathcal{GR},\mathcal{D}')$ 既非单射（它消去\SourcePage{199}$Vac'$ 在 $\bar U'$ 的中心 $\mathfrak{z}$ 中的零化子 $I$），也非满射（其像不含 $\operatorname{Der}O$）。我们把它分解为 $\bar U'\to U^\natural\to\Gamma(\mathcal{GR},\mathcal{D}')$，其中 $U^\natural$ 是在 $\bar U'/I\bar U'$ 上「添加」3.6.5 的代数胚 $I/I^2$ 而得（$\mathfrak{z}_{\mathfrak{g}}(O)=\mathfrak{z}/I\subset\bar U'/I\bar U'$ 与 $I/I^2$ 之间的交换关系来自 $I/I^2$ 上的代数胚结构，它们几乎是 Heisenberg 型的）。真空表示 $Vac'$ 作为 $U^\natural$-模不可约；$\Gamma(\mathcal{GR},I_\chi\lambda^{-1})$ 亦然，$\chi\in P_+({}^{L}G)$。
\end{enumerate}

\bdp{5.2.18} 以下是 5.2.16 的证明思路（我们假定 5.2.14）。置 $\mathfrak{z}:=\mathfrak{z}_{\mathfrak{g}}(O)$。考虑 $\mathfrak{z}$-模 $V^\chi_{\mathcal{L}\mathfrak{z}}:=\operatorname{Hom}_{\bar U'}(Vac',\allowbreak\Gamma(\mathcal{GR},I_\chi\lambda_{\mathcal{L}}^{-1}))$，于是 $\Gamma(\mathcal{GR},I_\chi\lambda_{\mathcal{L}}^{-1})=Vac'\otimes_{\mathfrak{z}}V^\chi_{\mathcal{L}\mathfrak{z}}$。某种 Tannakian 形式体系连同 Satake 等价，给出 $\operatorname{Spec}\mathfrak{z}$ 上典范的 ${}^{L}G$-torsor $\mathfrak{F}_{\mathfrak{z}}$，使得诸 $V^\chi_{\mathcal{L}\mathfrak{z}}$ 是 $V^\chi$ 的 $\mathfrak{F}$-扭。$\Gamma(\mathcal{GR},I_\chi\lambda_{\mathcal{L}}^{-1})$ 上的 $U^\natural$-模结构，定义了李代数胚 $I/I^2$ 在 $\mathfrak{F}_{\mathfrak{z}}$ 上的作用。一些额外的几何考虑在 $\mathfrak{F}_{\mathfrak{z}}$ 上定义了典范的 $B$-结构，它关于 $\operatorname{Der}O\subset I/I^2$ 的作用满足「oper」性质。于是 3.5、3.6 的结果给出典范等同 $(\operatorname{Spec}\mathfrak{z},\mathfrak{F}_{\mathfrak{z}})\simeq(\operatorname{Spec}A,\mathfrak{F}_A)$，其中 $A\simeq\mathfrak{z}$ 是 Feigin--Frenkel 同构。证毕。

\bdp{5.2.19} 我们需要它吗？？？
% [Tutor Note: 原书 p.199 此处标题为「DO WE NEED IT???」，是作者留在正文中的自问（三个问号照录），底本如此，保留原样。]

下面给出 $M$ 的一个直接构造，不用扭 $\mathcal{D}$-模。对 $x\in X$，考虑概形 $\operatorname{Bun}_{G,\bar x}$（见 2.3.1）。对 $\mathcal{L}\in Z\operatorname{tors}_\theta(X)$，记 $\lambda_{\mathcal{L},\bar x}$ 为线丛 $\lambda_{\mathcal{L}}$ 到 $\operatorname{Bun}_{G,\bar x}$ 的拉回。设 $\widetilde{\mathfrak{g}\otimes K_x}$ 是 2.5.1 中 $\mathfrak{g}\otimes K_x$ 的中心扩张。于是 $\mathfrak{g}\otimes K_x$ 在 $\operatorname{Bun}_{G,\bar x}$ 上的作用，典范地提升为 $\widetilde{\mathfrak{g}\otimes K_x}$ 在 $\lambda_{\mathcal{L},\bar x}$ 上的作用，其中 $\mathbf{1}\in\mathbb{C}$ 的作用为恒等（见 4.4.12）。记 $\operatorname{Bun}_{G,\mathcal{L},\bar x}$ 为 $\operatorname{Bun}_{G,\bar x}$ 上对应于 $\lambda_{\mathcal{L},\bar x}$ 的 $\mathbb{G}_m$-torsor 的全空间。我们有 Harish-Chandra 对 $(\widetilde{\mathfrak{g}\otimes K_x},\mathbb{G}_m\times G(O_x))$，其中 $\operatorname{Lie}\mathbb{G}_m=\mathbb{C}\subset\widetilde{\mathfrak{g}\otimes K_x}$。$\widetilde{\mathfrak{g}\otimes K_x}$ 在 $\operatorname{Bun}_{G,\mathcal{L},\bar x}$ 上的作用，按显然的方式延拓为该对的作用。

注意 $\operatorname{Bun}_G=\mathbb{G}_m\times G(O_x)\backslash\operatorname{Bun}_{G,\mathcal{L},\bar x}$。因此按 1.2.4 与 1.2.6，有函子 $\Delta_{\mathcal{L}}:\allowbreak\left(\widetilde{\mathfrak{g}\otimes K_x},\mathbb{G}_m\times G(O_x)\right)\allowbreak\operatorname{mod}\to\mathcal{M}^\ell(\operatorname{Bun}_G)$。

\SourcePage{200}

把投影 $\mathbb{G}_m\times G(O_x)\to\mathbb{G}_m$ 看作一个特征标；记 $Vac^\sim$ 为相应的诱导 Harish-Chandra 模。有
\begin{equation}\tag{258}
M_{\mathcal{L}}=\Delta_{\mathcal{L}}(Vac^\sim).
\end{equation}

下面确定 $M_{\mathcal{L}}$ 上的 $A_{{}^{L}\mathfrak{g}}(X)$-模结构。$\operatorname{End}(Vac^\sim)=\mathfrak{z}_{\mathfrak{g}}(O_x)$ 在 $\Delta_{\mathcal{L}}(Vac^\sim)$ 上的作用，经 Feigin--Frenkel 同构 $\varphi_{O_x}$（见 3.2.2），等同于一个 $A_{{}^{L}\mathfrak{g}}(O_x)$-作用。该作用经 $A_{{}^{L}\mathfrak{g}}(X)$ 这个商分解。

\bdsec{5.3. Satake 等价}

本节回顾基本事实与构造，并固定记号。细节与证明见 [MV]。[MV] 的作者用反常层；我们用 $\mathcal{D}$-模。

\bdp{5.3.1} 考虑仿射（或环路）Grassmann 流形 $\mathcal{GR}=G(K)/G(O)$（按惯例 $K=\mathbb{C}((t))$，$O=\mathbb{C}[[t]]$）；它是形式光滑、归纳射影的归纳概形（见 4.5.1）。它带有按 $G(O)$-轨道 $\operatorname{Orb}_\chi$（$\chi\in P_+({}^{L}G)$）的分层（见 4.5.8）。每个分层都是 $\operatorname{Aut}^0O$-不变的。

在 4.5.10 中我们引入了 $\mathcal{GR}$ 连通分支的奇偶性概念。按 4.5.11，
\begin{equation}\tag{259}
\begin{array}{c}
\text{$\mathcal{GR}$ 的偶（相应地奇）连通分支的所有分层，}\\[2pt]
\text{维数都是偶（相应地奇）的。}
\end{array}
\end{equation}

\bdp{5.3.2} \emph{引理.}
\begin{enumerate}[label=(\roman*),leftmargin=2.2em]
\item 每个分层 $\operatorname{Orb}_\chi$ 连通且单连通。
\item $\operatorname{Orb}_\chi$ 上任一光滑 $\mathcal{D}$-模是常值的。
\item $\operatorname{Orb}_\chi$ 只在偶次有上同调。
\end{enumerate}

\emph{证明.} 记 $\operatorname{Stab}_x$ 为 $x\in\mathcal{GR}$ 在 $G(O)$ 中的稳定化子。$\operatorname{Stab}_x$ 在 $G(O/tO)=G$ 中的像是抛物子群 $P_x$，而态射 $G(O)/\operatorname{Stab}_x\to G/P_x$ 是局部平凡纤维化，其纤维同构于仿射空间。于是 (i) 与 (iii) 显然。注意 $\overline{\operatorname{Orb}}_\chi$ 射影，且按 (259)，$\overline{\operatorname{Orb}}_\chi\setminus\operatorname{Orb}_\chi$ 的余维 $\ge 2$。故按
% 本单元止于 p.200 末句「So by」；下文（the projection formula …）见 p.201，由下一单元承接。
